\subsection{Bounded Sequences}

\begin{tcolorbox}[title=Problem 1, breakable]
    Show that if $(x_n)$ and $(y_n)$ are bounded sequences 
    and $c \in \mathbb{R}$, then the sequence $(cx_n)$ and $(x_n - y_n)$ are also bounded.
\end{tcolorbox}

\begin{proof}
    Notice $c = (c_n)$ where $c_n = c$ for all $n$.
    By Theorem 3.1.5, we have $|c_n| \le K = |c|$ and $|x_n| \le K'$
        so $|x_n c_n| \le K K'$.
    In addition $|y_n| \le K''$ so $|-y_n| \le K''$ for all $n$ and by the same Theorem we have 
    $|x_n - y_n| = |x_n + (-y_n)| \le K' + K''$.
\end{proof}

\begin{tcolorbox}[title=Problem 2, breakable]
    True or false (explain): If $(x_n)$ and $(x_n y_n)$ are bounded,
    then $(y_n)$ is also bounded.
\end{tcolorbox}

\begin{proof}
    False. Consider $x_n = 0$ and $y_n = 2^n$.
    Then $(x_n)$ and $(x_n y_n) = (0)$ are both bounded, but $(y_n)$ is not.
\end{proof}

\begin{tcolorbox}[title=Problem 3, breakable]
    Show that the sequence $(x_n)$ defined by  
    \[x_n = 1 + 1/2  + 1/3 + \cdots + 1/n,\]
    is unbounded.
\end{tcolorbox}

\begin{proof}
    We have $x_{2n} - x_n = \frac{1}{n+1} + \frac{1}{n+2} + \cdots + \frac{1}{2n} \ge \frac{n}{2n} = \frac{1}{2}$ for all $n$.
    Suppose $(x_n)$ is bounded by $M \ge 0$. By the Archimedean property, there exists $k \in \mathbb{N}$ such that $k \cdot \frac{1}{2} > M$.
    But then $x_{2^k} \ge x_1 + k \cdot \frac{1}{2} > M$, a contradiction.
\end{proof}

\begin{tcolorbox}[title=Problem 6, breakable]
    Show that the sequence $(x_n)$ defined by
    \[x_n = 1 + \frac{1}{2!} + \frac{1}{3!} + \cdots + \frac{1}{n!}\]
    is bounded above by $2$.
\end{tcolorbox}

\begin{proof}
    Notice $n! \ge 2^{n-1}$ for all $n$, so $\frac{1}{n!} \le \frac{1}{2^{n-1}}$. Thus
    \[x_n \le \sum_{k=1}^{n} \frac{1}{2^{k-1}} = \frac{1 - (1/2)^n}{1 - 1/2} 
    = 2\left(1 - \frac{1}{2^n}\right) < 2.\]
\end{proof}

\subsection{Ultimately, Frequently}

\newpage
\begin{tcolorbox}[title=Problem 1, breakable]
    For the sequence $x_n = n$, $x_n$ is both frequently even and frequently odd;  
        does this conflict with 3.2.4.
\end{tcolorbox}

\begin{proof}
    No.
\end{proof}

\begin{tcolorbox}[title=Problem 3, breakable]
    Let $(A_n)$ be a sequence of subsets of a set $X$. Define
    \begin{align*}
        \limsup A_n &= \{x \in X : x \in A_n \text{ frequently}\}, \\
        \liminf A_n &= \{x \in X : x \in A_n \text{ ultimately}\}.
    \end{align*}
    Prove:
    \begin{enumerate}[(i)]
        \item $\liminf A_n \subseteq \limsup A_n$.
        \item $\liminf A_n = \bigcup_{n=1}^{\infty} \bigcap_{k=n}^{\infty} A_k$.
        \item $(\limsup A_n)' = \liminf A_n'$ \quad ($'$ means complement).
        \item $\limsup A_n = \bigcap_{n=1}^{\infty} \bigcup_{k=n}^{\infty} A_k$.
    \end{enumerate}
\end{tcolorbox}

\begin{proof}
     Let $x \in \liminf A_n$.
     Now choose $N$ such that $x_n \in \liminf A_N$ which can be done 
        since $  \liminf A_n$ ultimately appears in the indices.
    Then since $x_n$ appears frequently in $\liminf A_n$ it must appear for some $N' > N$
        for $\liminf A_n$.
    That is, $x_n \in \liminf A_N$.
\end{proof}

\begin{proof}
    Let $x \in \liminf A_n$.
    Then $x \in A_n$ ultimately, meaning there exists $N$ such that $x \in A_n$ for all $n \ge N$.
    That is, $x \in \bigcap_{k=N}^{\infty} A_k$, so $x \in \bigcup_{n=1}^{\infty} \bigcap_{k=n}^{\infty} A_k$.
    Conversely, if $x \in \bigcup_{n=1}^{\infty} \bigcap_{k=n}^{\infty} A_k$, then there exists $N$ such that
    $x \in \bigcap_{k=N}^{\infty} A_k$, meaning $x \in A_n$ for all $n \ge N$, so $x \in A_n$ ultimately,
    thus $x \in \liminf A_n$.
\end{proof}

\begin{proof}
    Let $x \in (\limsup A_n)'$.
    Then $x \notin \limsup A_n$, so $x \in A_n$ only finitely many times, so \ $x \notin A_n$ ultimately.
    That is, there exists $N$ such that $x \notin A_n$ for all $n \ge N$, meaning $x \in A_n'$ for all $n \ge N$.
    Thus $x \in A_n'$ ultimately, so $x \in \liminf A_n'$.
    The reverse inclusion is identical with all implications reversed.
\end{proof}

\begin{proof}
    Let $x \in \limsup A_n$.
    Then $x \in A_n$ frequently, meaning for every $N$ there exists $k \ge N$ such that $x \in A_k$.
    Thus $x \in \bigcup_{k=N}^{\infty} A_k$ for every $N$, so $x \in \bigcap_{n=1}^{\infty} \bigcup_{k=n}^{\infty} A_k$.
    Conversely, if $x \in \bigcap_{n=1}^{\infty} \bigcup_{k=n}^{\infty} A_k$, then for every $N$ we have
    $x \in \bigcup_{k=N}^{\infty} A_k$, so there exists $k \ge N$ with $x \in A_k$, meaning $x \in A_n$ frequently,
    thus $x \in \limsup A_n$.
\end{proof}

\subsection{Null Sequences}

\begin{tcolorbox}[title=Problem 1, breakable]
    Prove: If $(x_n)$ and $(y_n)$ are null then so is $(x_n y_n)$.
\end{tcolorbox}
\begin{tcolorbox}[title=Problem 2, breakable]
    Prove that $(x_n)$ is null if and only if $(x_n^2)$ is null. (Generalization?)
\end{tcolorbox}
\begin{tcolorbox}[title=Problem 3, breakable]
    True or false (explain): If $(x_n)$ and $(y_n)$ are sequences in $\mathbb{R}$ such
    that $(x_n y_n)$ is null, then either $(x_n)$ or $(y_n)$ is null.
\end{tcolorbox}
\begin{tcolorbox}[title=Problem 5, breakable]
    Show that a sequence $(x_n)$ in $\mathbb{R}$ is \textit{not} null if and only if there
    exists a positive number $\epsilon$ such that $|x_n| \ge \epsilon$ frequently.
\end{tcolorbox}
\begin{tcolorbox}[title=Problem 11, breakable]
    If $p : \mathbb{R} \to \mathbb{R}$ is a polynomial function without constant term and
    $(x_n)$ is a null sequence, then $(p(x_n))$ is null.
\end{tcolorbox}